% --- Άσκηση 35042 (35042.tex) ---
\Askhsh{35042}{2}
\begin{ylh}{2}
    \item 3.3. Εξισώσεις 2ου Βαθμού
    \item 5.1. Ακολουθίες
    \item 5.3. Γεωμετρική πρόοδος
\end{ylh}

\begin{alist}
\item Να βρείτε, για ποιες τιμές του $x$, οι αριθμοί $x + 4,\text{ }\ 2 - x,\text{ }\ 6 - x$ με τη σειρά που δίνονται είναι διαδοχικοί όροι γεωμετρικής προόδου.
\monades{13}
\item Αν $x = 5$ και ο $6 - x$ είναι ο τέταρτος όρος της παραπάνω γεωμετρικής προόδου, να βρείτε:
\begin{rlist}
\item
  το λόγο $\lambda$ της γεωμετρικής προόδου.
\monades{6}
\item
  τον πρώτο όρο $a_{1}$ της προόδου.
\monades{6}
\end{rlist}
\end{alist}
